\section{Related Work}
\label{sec:related}

\paragraph{Formal Methods and Verified Systems.}
The verification of software against formal specifications is a
mature field. TLA+~\cite{lamport2002} and its model
checker TLC verify safety and liveness properties of distributed
systems; Alloy~\cite{jackson2012} expresses relational invariants
and finds counter-examples by bounded model checking;
dependently typed languages such as Coq~\cite{bertot2004} and
Agda enable machine-checked proofs of program properties. These
tools verify that a \emph{given specification} satisfies logical
invariants. Constitutional Code differs in scope: it does not verify
that a particular implementation is logically correct, but that the
\emph{architectural separation of capabilities between principals}
satisfies a political theory requirement. TLA+ can verify Executive
Containment (Theorem~\ref{thm:exec}) for a specific operator binary;
it cannot, within its specification language, certify that the
resulting institution satisfies Legislative Independence as a
political property, because the concept of ``a principal who may
not redefine the rules'' has no primitive representation in TLA+.
Our contribution is the mapping between political theory and
cryptographic architecture.

\paragraph{Algorithmic Accountability and Transparency.}
Substantial work documents the harms of opaque algorithmic
systems~\cite{pasquale2015,obermeyer2019dissecting} and proposes
normative frameworks for accountability~\cite{doshi2017,
wachter2017}. The ``right to explanation'' literature traces the
legal basis for algorithmic transparency under
GDPR~\cite{wachter2017}. Kearns and Roth~\cite{kearns2019ethical}
explore the tension between privacy and fairness in machine
learning. These works define normative goals but stop short of
proposing enforcement mechanisms that are resistant to the
operator's self-interest. Constitutional Code provides the missing
enforcement layer.

\paragraph{Smart Contracts and Algorithmic Law.}
Szabo~\cite{szabo1997} proposed the concept of self-enforcing
contracts implemented as code. Blockchain-based smart contracts
~\cite{buterin2014} realise this vision for financial transactions.
However, smart contract systems address \emph{consensus}
(agreement among nodes on a shared state) rather than
\emph{accountability} (verifying that a specific operator's
decisions comply with individual-rights obligations). They also
lack privacy-preserving auditability: the full transaction history
is public, making them unsuitable for high-stakes AI decisions
involving sensitive personal data. Constitutional Code addresses
both gaps through the combination of linear types and ZK proofs.

\paragraph{Zero-Knowledge Proofs in Governance.}
Recent work applies ZK proofs to privacy-preserving auditing in
specific domains: tax compliance~\cite{bunz2018bulletproofs},
voting~\cite{benaloh1994receipt}, and financial
regulation~\cite{zkledger2018}. Paper~3 of this series~\cite{btv-ar2026}
extends this to privacy-preserving algorithmic accountability at
the batch level. Constitutional Code situates Paper~3's construction
within a broader institutional framework, showing that ZK
auditability is not merely a privacy tool but a constitutional
necessity for Judicial Independence.

\paragraph{Political Theory and Computation.}
Lessig's \emph{Code and Other Laws of Cyberspace}~\cite{lessig1999}
argued that code \emph{is} law---that the architecture of information
systems functions as regulation. Constitutional Code operationalises
this insight in the reverse direction: \emph{law can be architecture}.
Montesquieu's separation of powers is not merely encoded into code as
a metaphor; it is proven as a theorem about capability disjointness
in a formal type system.
